% GENERATED FILE. Edit canonical lessons, then run npm run generate:books.
% canonical-source-hash: fnv1a64:858cd36054cc183d
% canonical-lessons: GU-C413-taraph, GU-C413-marg, GU-C413-par, GU-C413-dorvu, GU-C413-cidhvu

\chapter{Which Way?}
\label{ch:gu-which-way}
\hlchaptermodality{413}

\begin{chapteropening}
\textbf{By the end of this chapter:} \emph{I can ask the way and follow someone who points it out.}

Complete the last lesson of chapter 413: I can ask the way and follow someone who points it out.
\end{chapteropening}

\section[\texorpdfstring{\gu{તરફ} (taraph)}{તરફ (taraph)}]{\gu{તરફ} (taraph) — towards; in the direction of}

\label{lesson:GU-C413-taraph}



\begin{quote}
Before the new one: say the Gujarati for to drive in, then the Gujarati for to spend (time).
\end{quote}

\subsection*{You'll want to know: \gu{તરફ}}
\textbf{\gu{તરફ}} — \emph{taraph} — \textquotedblleft{}towards; in the direction of\textquotedblright{}.

\begin{quote}
\gu{અમે} \gu{નદી} \gu{તરફ} \gu{ચાલ્યા}.
\end{quote}

\begin{quote}
\emph{ame nadī taraph cālyā.}
\end{quote}

\begin{quote}
We walked towards the river.
\end{quote}

\begin{cousinweb}
\textbf{\gu{તરફ}} (\emph{taraph}) is usually traced to Arabic \emph{ṭaraf}, \textquotedblleft{}a side, an edge\textquotedblright{}, which came through Persian; after a noun it means towards it: \gu{ઘર} \gu{તરફ}, towards home.
\end{cousinweb}

\begin{grammarlens}[title={What you've built}]
One more word for getting around.
\end{grammarlens}

\subsection*{Your turn}
\begin{itemize}
\raggedright
  \item \emph{Say it:} \emph{taraph}
  \item \emph{Say it:} \emph{taraph}, once more
  \item \emph{Say it:} \emph{ame nadī taraph cālyā.} — We walked towards the river.
  \item \emph{Recall:} say the Gujarati for to drive in, then the Gujarati for to spend (time), then say \emph{taraph} again
\end{itemize}

\begin{tcolorbox}[breakable,colback=teal!4,colframe=teal!35!black,title={Before you move on}]
What does \gu{તરફ} mean? (\textquotedblleft{}Towards; in the direction of\textquotedblright{}.) Say it once more.
\end{tcolorbox}
\section[\texorpdfstring{\gu{માર્ગ} (mārg)}{માર્ગ (mārg)}]{\gu{માર્ગ} (mārg) — a route; a way (formal)}

\label{lesson:GU-C413-marg}



\begin{quote}
Before the new one: say the Gujarati for to spend (time), then the Gujarati for towards.
\end{quote}

\subsection*{You'll want to know: \gu{માર્ગ}}
\textbf{\gu{માર્ગ}} — \emph{mārg} — \textquotedblleft{}a route; a way (formal)\textquotedblright{}.

\begin{quote}
\gu{શહેરમાં} \gu{જવાનો} \gu{સૌથી} \gu{ટૂંકો} \gu{માર્ગ} \gu{કયો} \gu{છે}?
\end{quote}

\begin{quote}
\emph{śahermā̃ javāno sauthī ṭū̃ko mārg kayo che?}
\end{quote}

\begin{quote}
What is the shortest route into the city?
\end{quote}

A masculine noun.

\begin{cousinweb}
\textbf{\gu{માર્ગ}} (\emph{mārg}) is a tatsama, Sanskrit \emph{mārga}, first the track of a wild animal and then any road or way (Monier-Williams); signs and books use it where speech says \gu{રસ્તો}.
\end{cousinweb}

\begin{grammarlens}[title={What you've built}]
One more word for getting around.
\end{grammarlens}

\subsection*{Your turn}
\begin{itemize}
\raggedright
  \item \emph{Say it:} \emph{mārg}
  \item \emph{Say it:} \emph{mārg}, once more
  \item \emph{Say it:} \emph{śahermā̃ javāno sauthī ṭū̃ko mārg kayo che?} — What is the shortest route into the city?
  \item \emph{Recall:} say the Gujarati for to spend (time), then the Gujarati for towards, then say \emph{mārg} again
\end{itemize}

\begin{tcolorbox}[breakable,colback=teal!4,colframe=teal!35!black,title={Before you move on}]
What does \gu{માર્ગ} mean? (\textquotedblleft{}A route; a way (formal)\textquotedblright{}.) Say it once more.
\end{tcolorbox}
\section[\texorpdfstring{\gu{પાર} (pār)}{પાર (pār)}]{\gu{પાર} (pār) — the far side (of a river or a road)}

\label{lesson:GU-C413-par}



\begin{quote}
Before the new one: say the Gujarati for towards, then the Gujarati for a route.
\end{quote}

\subsection*{You'll want to know: \gu{પાર}}
\textbf{\gu{પાર}} — \emph{pār} — \textquotedblleft{}the far side (of a river or a road)\textquotedblright{}.

\begin{quote}
\gu{નદી} \gu{પાર} \gu{એક} \gu{નાનું} \gu{ગામ} \gu{છે}.
\end{quote}

\begin{quote}
\emph{nadī pār ek nānũ gām che.}
\end{quote}

\begin{quote}
There is a small village on the far side of the river.
\end{quote}

A masculine noun.

\begin{cousinweb}
\textbf{\gu{પાર}} (\emph{pār}) is a tadbhava of Sanskrit \emph{pāra}, \textquotedblleft{}the further bank, the furthest end\textquotedblright{}; CDIAL lists the Gujarati word (CDIAL 8100).
\end{cousinweb}

\begin{grammarlens}[title={What you've built}]
One more word for getting around.
\end{grammarlens}

\subsection*{Your turn}
\begin{itemize}
\raggedright
  \item \emph{Say it:} \emph{pār}
  \item \emph{Say it:} \emph{pār}, once more
  \item \emph{Say it:} \emph{nadī pār ek nānũ gām che.} — There is a small village on the far side of the river.
  \item \emph{Recall:} say the Gujarati for towards, then the Gujarati for a route, then say \emph{pār} again
\end{itemize}

\begin{tcolorbox}[breakable,colback=teal!4,colframe=teal!35!black,title={Before you move on}]
What does \gu{પાર} mean? (\textquotedblleft{}The far side (of a river or a road)\textquotedblright{}.) Say it once more.
\end{tcolorbox}
\section[\texorpdfstring{\gu{દોરવું} (dorvũ)}{દોરવું (dorvũ)}]{\gu{દોરવું} (dorvũ) — to lead (someone along); to guide}

\label{lesson:GU-C413-dorvu}



\begin{quote}
Before the new one: say the Gujarati for a route, then the Gujarati for the far side.
\end{quote}

\subsection*{You'll want to know: \gu{દોરવું}}
\textbf{\gu{દોરવું}} — \emph{dorvũ} — \textquotedblleft{}to lead (someone along); to guide\textquotedblright{}.

\begin{quote}
\gu{ભોમિયો} \gu{અમને} \gu{પર્વત} \gu{પર} \gu{દોરી} \gu{ગયો}.
\end{quote}

\begin{quote}
\emph{bhomiyo amne parvat par dorī gayo.}
\end{quote}

\begin{quote}
The guide led us up the mountain.
\end{quote}

Past (he …): \gu{તે} \gu{બાળકને} \gu{ઘરે} \gu{દોરી} \gu{ગયો}, \emph{te bāḷakne ghare dorī gayo}.

\begin{cousinweb}
\textbf{\gu{દોરવું}} (\emph{dorvũ}) goes back to \emph{*dōr}, \textquotedblleft{}to look at or for\textquotedblright{}, a form CDIAL reconstructs; CDIAL lists the Gujarati verb, \textquotedblleft{}to direct\textquotedblright{} (CDIAL 6581). It also means to draw a line, a sense this lesson leaves for later.
\end{cousinweb}

\begin{grammarlens}[title={What you've built}]
One more word for getting around.
\end{grammarlens}

\subsection*{Your turn}
\begin{itemize}
\raggedright
  \item \emph{Say it:} \emph{dorvũ}
  \item \emph{Say it:} \emph{dorvũ}, once more
  \item \emph{Say it:} \emph{bhomiyo amne parvat par dorī gayo.} — The guide led us up the mountain.
  \item \emph{Recall:} say the Gujarati for a route, then the Gujarati for the far side, then say \emph{dorvũ} again
\end{itemize}

\begin{tcolorbox}[breakable,colback=teal!4,colframe=teal!35!black,title={Before you move on}]
What does \gu{દોરવું} mean? (\textquotedblleft{}To lead (someone along); to guide\textquotedblright{}.) Say it once more.
\end{tcolorbox}
\section[\texorpdfstring{\gu{ચીંધવું} (cī̃dhvũ)}{ચીંધવું (cī̃dhvũ)}]{\gu{ચીંધવું} (cī̃dhvũ) — to point out (with a finger)}

\label{lesson:GU-C413-cidhvu}



\begin{quote}
Before the new one: say the Gujarati for the far side, then the Gujarati for to lead.
\end{quote}

\subsection*{You'll want to know: \gu{ચીંધવું}}
\textbf{\gu{ચીંધવું}} — \emph{cī̃dhvũ} — \textquotedblleft{}to point out (with a finger)\textquotedblright{}.

\begin{quote}
\gu{એક} \gu{માણસે} \gu{અમને} \gu{સ્ટેશનનો} \gu{રસ્તો} \gu{ચીંધ્યો}.
\end{quote}

\begin{quote}
\emph{ek māṇse amne sṭeśanno rasto cī̃dhyo.}
\end{quote}

\begin{quote}
A man pointed out the way to the station.
\end{quote}

Past, with \gu{તેણે} (he or she …): \gu{તેણે} \gu{રસ્તો} \gu{ચીંધ્યો}, \emph{teṇe rasto cī̃dhyo}.

\begin{cousinweb}
\textbf{\gu{ચીંધવું}} (\emph{cī̃dhvũ}) No dictionary we use records an older form, so its history is not settled.
\end{cousinweb}

\begin{grammarlens}[title={What you've built}]
One more word for getting around.
\end{grammarlens}

\subsection*{Your turn}
\begin{itemize}
\raggedright
  \item \emph{Say it:} \emph{cī̃dhvũ}
  \item \emph{Say it:} \emph{cī̃dhvũ}, once more
  \item \emph{Say it:} \emph{ek māṇse amne sṭeśanno rasto cī̃dhyo.} — A man pointed out the way to the station.
  \item \emph{Recall:} say the Gujarati for the far side, then the Gujarati for to lead, then say \emph{cī̃dhvũ} again
\end{itemize}

\begin{tcolorbox}[breakable,colback=teal!4,colframe=teal!35!black,title={Before you move on}]
What does \gu{ચીંધવું} mean? (\textquotedblleft{}To point out (with a finger)\textquotedblright{}.) Say it once more.
\end{tcolorbox}
